\begin{center}
	\LARGE \bf {Appendix}
\end{center}

\etocdepthtag.toc{mtappendix}
\etocsettagdepth{mtchapter}{none}
\etocsettagdepth{mtappendix}{subsection}
\tableofcontents
\newpage

\section{Broad Impact}
This work explores the integration of KGs with LLMs for the task of KGQA, with a focus on improving retrieval quality and generalization through graph-based retrievers. By grounding LLMs in structured relational knowledge and employing a more interpretable and generalizable retrieval model, our proposed framework aims to significantly mitigate hallucination and improve factual consistency.

The potential societal benefits of this research are considerable. More faithful and robust KG-augmented language systems could enhance the reliability of AI assistants in high-stakes domains such as scientific research, healthcare, and legal reasoning—where factual accuracy and interpretability are critical. In particular, reducing hallucinations in LLM outputs can support safer deployment of AI systems in real-world applications.
Moreover, our method is efficient at inference stage, and does not rely on large-scale supervised training or fine-tuning of LLMs, making it more accessible and sustainable, especially in resource-constrained environments. Additionally, by supporting smaller LLMs with strong retrieval capabilities, our work helps democratize access to high-performance question answering systems without requiring access to proprietary or computationally expensive language models for knowledge retrieval.


\section{Related Work}
\label{related_work}
\tb{Retrieve-then-reasoning paradigm}.
In KG-based RAG, a substantial number of methods adopt the \textit{retrieve-then-reasoning} paradigm~\cite{li2023graph,kim-etal-2023-kg,liu2025dualreasoninggnnllmcollaborative,wu2023retrieverewriteanswerkgtotextenhancedllms,wen2023mindmap,mavromatis2024gnn,li2024simple,luo2024rog,mavromatis2022rearev}, wherein a retriever first extracts relevant triples from the knowledge graph, followed by a (LLM-based) that generates the final answer based on the retrieved information. The retriever can be broadly categorized into LLM-based and graph-based approaches.
LLM-based retrievers benefit from large model capacity and strong semantic understanding. However, they also suffer from hallucination issues, and high computational cost and latency. To address this, recent work has proposed lightweight GNN-based graph retrievers~\cite{mavromatis2022rearev,li2024simple,zhang2022subgraph,mavromatis2022rearev} that operate directly on the KG structure. These methods have achieved superior performance on KGQA benchmarks, benefiting from the strong reasoning and denoising capabilities of downstream LLM-based reasoners. While graph-based retrievers offer substantial computational advantages and mitigating hallucinations, they still face generalization challenges~\cite{li2024simple}. In this work, we aim to retain the efficiency of graph-based retrievers while  enhancing their expressivity and generalization ability through a model-agnostic design tailored for characteristics of knowledge graphs and KGQA task.

\tb{KG-based agentic RAG}. 
Another line of research leverages LLMs as agents that iteratively explore the knowledge graph to retrieve relevant information~\cite{gao2024two,wang2024knowledge,jiang2024kg,sunthink,chen2024plan,ma2024think,jin2024graph}. In this setting, the agent integrates both retrieval and reasoning capabilities, enabling more adaptive knowledge access. While this approach has demonstrated effectiveness in identifying relevant triples, the iterative exploration process incurs substantial latency, as well as computational costs due to repeated LLM calls. In contrast, our method adopts a lightweight graph-based retriever to balance efficiency and effectiveness for KGQA.

\section{Complexity Analysis}
\label{app:complexity}
\tb{Preprocessing.} Given a graph $\mcal{G} = (\mcal{V}, \mcal{E})$ with $|\mcal{V}|$ nodes and $|\mcal{E}|$ edges,  The time complexity of this transformation is $\mathcal{O}(|\mcal{E}| d_{\max})$, where $d_{\max}$ is the maximum node degree in $\mcal{G}$. The space complexity is $\mathcal{O}(|\mcal{E}| + |\mcal{E}'|)$, where $|\mcal{E}'|$ denotes the number of edges in the resulting line graph $\mcal{G}'$, typically on the order of $\mathcal{O}(|E| d_{avg})$, with $d_{avg}$ as the average node degree.

\tb{Model training and inference.} For both training and inference, with a $K$-layer GCN operating on the line graph $\mcal{G}_q'$, the time complexity is $\mathcal{O}(K|\mcal{E}'| F)$, where $F$ is the dimensionality of node features. The space complexity is $\mathcal{O}(|\mcal{V}'|F + |\mcal{E}'|)$, where $|V'|$ is the number of triples (i.e., nodes in the line graph). During model training and inference, it does not involve any LLM call for the retriever. 

\section{More Discussions on Line Graph Transformation}
\label{discuss_graph_trasnform}
In this section, we discuss the limitations of additional message-passing GNNs when applied to the original knowledge graph, and elaborate on how these limitations can be addressed through line graph transformation.

\tb{GraphSage~\cite{hamilton2017inductive}}. The upadte function of GraphSage is shown in Eqn.~\ref{eq:sage_edge_concat}. 
\begin{equation}\label{eq:sage_edge_concat}
  \mathbf m_{ij}^{(k)}
    = W_{\text{msg}}^{(k)}
      \bigl[\,
        \mathbf h_j^{(k-1)} \,\Vert\, \mathbf r_{ij}
      \bigr],
  \qquad
  \mathbf h_i^{(k)}
    = \sigma\!\Bigl(
        W_{\text{self}}^{(k)}\mathbf h_i^{(k-1)}
        \;+\;
        \mathrm{MEAN}\!\bigl\{
          \mathbf m_{ij}^{(k)}
        \bigr\}
      \Bigr).
\end{equation}
As seen, only the entity states $\mathbf h_i^{(k)}$ are upadted, and relation embeddings $\mathbf r_{ij}$ remain static, so the model cannot refine relation semantics nor capture cross-triple interactions. 
After converting to the line graph $\mathcal G'_q$, the propagation is performed between triples, enabling explicit inter-triple reasoning.
\begin{equation}\label{eq:sage_line_concat}
  \mathbf h_u^{(k)}
    = \sigma\!\Bigl(
        W_{\text{self}}^{(k)}
        \,\phi\bigl(\mathbf e_u,\mathbf r_u,\mathbf e_u'\bigr)
        \;+\;
        W_{\text{nbr}}^{(k)}
        \frac{1}{|\mathcal N_\ell(u)|}
        \sum_{v\in\mathcal N_\ell(u)}
          \phi\bigl(\mathbf e_v,\mathbf r_v,\mathbf e_v'\bigr)
      \Bigr),
\end{equation}

\tb{Graph Attention Network (GAT)~\cite{velivckovic2017graph}}. The upadte function of GAT is shown in Eqn.~\ref{eq:gat_edge_concat}. Similarly $\mathbf r_{ij}$ only modulates one-hop attention, therefore the relational embedding is not updated, only the entity embedding is learned. 
\begin{align}\label{eq:gat_edge_concat}
  \alpha_{ij}^{(k)}
    &= \mathrm{softmax}_{j}\!
       \Bigl(
         a^\top
         \bigl[
           W \mathbf h_i^{(k-1)} \,\Vert\,
           W \mathbf h_j^{(k-1)} \,\Vert\,
           W_e \mathbf r_{ij}
         \bigr]
       \Bigr), \nonumber\\
  \mathbf h_i^{(k)}
    &= \sigma\!\Bigl(
        \sum_{j\in\mathcal N(i)}
          \alpha_{ij}^{(k)}\;
          V\bigl[\,
            \mathbf h_j^{(k-1)} \,\Vert\, \mathbf r_{ij}
          \bigr]
      \Bigr).
\end{align}
After the line-graph transform, Attention is now computed between triples, so both intra-triple composition and inter-triple dependencies influence edge-aware attention.

\begin{equation}
  \alpha_{ij}
  \;=\;
  \operatorname{softmax}_{h_j}
  \!\bigl(
    a^\top\!
    \bigl[
      W\phi(\mathbf e_i,\mathbf r_i,\mathbf e_i')
      \;\Vert\;
      W\phi(\mathbf e_j,\mathbf r_j,\mathbf e_j')
    \bigr]
  \bigr),
  \qquad
  \rvh_{i}^{(k)}
  \;=\;
  \sigma\!\Bigl(
    \sum_{h_j\in\mathcal N_\ell(h_i)}
      \alpha_{ij}\,
      W\phi(\mathbf e_j,\mathbf r_j,\mathbf e_j')
  \Bigr).
\end{equation}
Similarly, for all message-passing GNNs that follow the general \textit{message–aggregation–update} paradigm (Eq.~\ref{eq:general_gnn}), the relational structure in the original knowledge (sub)graph cannot be fully exploited. In contrast, the line graph transformation offers a model-agnostic solution that enables richer modeling of both intra- and inter-triple interactions, facilitating more expressive relational reasoning. This transformation is broadly applicable across a wide range of GNN architectures.
\begin{equation}
\begin{aligned}
\label{eq:general_gnn}
% ----------------------------------------------------
% 1)  Message construction
\mathbf m_{ij}^{(k)}
  &= \mathcal M^{(k)}
     \!\bigl(
       \mathbf h_i^{(k-1)},\;
       \mathbf h_j^{(k-1)},\;
       \mathbf r_{ij}
     \bigr),
     &&\forall\, j\!\in\!\mathcal N(i),
\\[4pt]
% ----------------------------------------------------
% 2)  Neighbour aggregation
\mathbf a_i^{(k)}
  &= \mathcal A^{(k)}
     \!\Bigl(
       \bigl\{\mathbf m_{ij}^{(k)} : j\!\in\!\mathcal N(i)\bigr\}
     \Bigr),
\\[4pt]
% ----------------------------------------------------
% 3)  Node/edge update
\mathbf h_i^{(k)}
  &= \mathcal U^{(k)}
     \!\bigl(
       \mathbf h_i^{(k-1)},\;
       \mathbf a_i^{(k)}
     \bigr),
\end{aligned}
\end{equation}


\section{Algorithmic Pseudocode}
\label{pseudocode}
The overall preprocessing and training procedure is illustrated in Algorithm~\ref{alg:proj}.

\begin{algorithm}[t]
\caption{Overall Procedure of \proj}
\label{alg:proj}
\begin{algorithmic}[1]
\Require Training set $\mathcal{D}=\{(q,e_q,e_a,\mathcal{G}_q)\}$, epochs $E$, learning rate $\eta$, hyper‑parameters $\lambda_q,\lambda_{path}$
\Ensure  Optimized retriever parameters $\theta=\{\overrightarrow{\theta},\overleftarrow{\theta}\}$ and $\phi$
%
\Statex\textbf{Preprocessing}
\ForAll{$(q,e_q,e_a,\mathcal{G}_q)\in\mathcal{D}$}
    \State Construct line graph $\mathcal{G}'_q \gets \textsc{LineGraph}(\mathcal{G}_q)$
    \State Generate candidate paths $\mathcal{P}_{cand} \gets \textsc{GenPaths}(\mathcal{G}_q,d_{\min},d_{\min}{+}2)$
    \State Rationalize labels $\mathcal{Y}_q \gets \gamma(q,\mathcal{P}_{cand})$ \Comment{LLM call}
    \State Relation targeting $\mathcal{R}_* \gets \gamma(q,\textsc{Relations}(\mathcal{G}_q))$ \Comment{LLM call}
\EndFor
%
\Statex\textbf{Initialization}
\State Initialize bi‑directional GCN encoders $f_{\overrightarrow{\theta}},\,f_{\overleftarrow{\theta}}$ and STOP MLP $g_\phi$
%
\Statex\textbf{Training}
\For{$e=1$ \textbf{to} $E$}
    \ForAll{minibatch $\mathcal{B}\subset\mathcal{D}$}
        \ForAll{$(q,\mathcal{G}'_q,\mathcal{Y}_q,\mathcal{R}_*)\in\mathcal{B}$}
            \State Encode nodes $\mathbf{z}_i \gets 
            \tfrac12\!\bigl(f_{\overrightarrow{\theta}}(v_i)+
            f_{\overleftarrow{\theta}}(v_i)\bigr)$
            \State Build $\mathcal{V}_{pos},\mathcal{V}_{neg}$ using $\mathcal{R}_*$   
            \State $\mathcal{L}_q \gets \textsc{NegSampleLoss}(\theta,\phi;\mathcal{V}_{pos},\mathcal{V}_{neg})$ \Comment{Eq.~\eqref{eq:L_q_neg_sampling}}
            \State $\mathcal{L}_{path} \gets \textsc{PathLoss}(\theta,\phi;q,\mathcal{Y}_q)$  \Comment{Eq.~\eqref{path_loss}}
        \EndFor
        \State $\mathcal{L} \gets \lambda_q\mathcal{L}_q + \lambda_{path}\mathcal{L}_{path}$ \Comment{$\lambda_q$, $\lambda_{path}$ default to $1.0$}
        \State Update $\theta,\phi$ via Adam with step size $\eta$
    \EndFor
\EndFor
%
\end{algorithmic}
\end{algorithm}


\section{Theoretical Properties of the Directed Line Graph}
Let $\mathcal{G}=(\mcal{V},\mcal{E})$ be a finite directed graph.
Its directed line graph is denoted by $l(\mathcal{G})=(\mcal{V}_{\ell},\mcal{E}_{\ell})$,
where every node $x\in\mcal{V}_{\ell}$ corresponds to a directed edge
$e_x=(u,v)\in\mcal{E}$, and there is an edge $(x,y)\in\mcal{E}_{\ell}$
iff the head of $e_x$ equals the tail of $e_y$.

\vspace{4pt}
\begin{proposition}[Bijective mapping of directed paths]
\label{prop:path-bijection}
Let
\[
P \;=\;
\bigl(
  v_0 \xrightarrow{\,e_1\,} v_1 \xrightarrow{\,e_2\,} \dots
  \xrightarrow{\,e_k\,} v_k
\bigr),\quad k\ge 1,
\]
be a directed path of length $k$ in $\mathcal{G}$, where
$e_i=(v_{\,i-1},v_i)\in\mcal{E}$.
Define
\[
\Phi(P)
\;=\;
\bigl(e_1 \xrightarrow{} e_2 \xrightarrow{}\dots\xrightarrow{} e_k\bigr),
\]
regarding each $e_i$ as its corresponding node in $l(\mathcal{G})$.
Then $\Phi$ is a bijection between the set of directed paths of length $k$ in $\mathcal{G}$, and
the set of directed paths of length $k-1$ in $l(\mathcal{G})$.
\end{proposition}

\begin{proof}
\emph{Well‑definedness.}  
For consecutive edges $e_i=(v_{i-1},v_i)$ and $e_{i+1}=(v_i,v_{i+1})$
the shared endpoint is $v_i$ with consistent orientation, hence
$(e_i,e_{i+1})\in\mcal{E}_{\ell}$.
Thus $\Phi(P)$ is a valid path of length $k-1$ in~$l(\mathcal{G})$.

\emph{Injectivity.}  
If two original paths $P\neq P'$ differ at position $j$, then
$\Phi(P)$ and $\Phi(P')$ differ at node~$j$, so $\Phi(P)\neq\Phi(P')$.

\emph{Surjectivity.}  
Take any path $(x_1\!\to x_2\!\to\dots\to x_k)$ in $l(\mathcal{G})$ and
let $e_i$ be the edge in $\mcal{E}$ represented by $x_i$.
Because $(x_i,x_{i+1})\in\mcal{E}_{\ell}$, the head of $e_i$ is the tail
of $e_{i+1}$.  Hence $(e_1,\dots,e_k)$ forms a length‑$k$ path in
$\mathcal{G}$ whose image under~$\Phi$ is the given path.
\end{proof}

\vspace{4pt}
\begin{proposition}[Path‑length reduction]
\label{prop:length-reduction}
Let $u,v\in\mcal{V}$ be connected in $\mathcal{G}$ by a shortest
directed path of length $d\ge 1$,
\(
u=v_0\!\xrightarrow{e_1}\!\dots\xrightarrow{e_d} v=v_d.
\)
Let $x_u$ and $x_v$ be the nodes of $\mcal{V}_{\ell}$ corresponding to
the first edge $e_1$ and the last edge $e_d$, respectively.
Then the distance between $x_u$ and $x_v$ in $l(\mathcal{G})$ is
exactly $d-1$, and no shorter path exists in $l(\mathcal{G})$.
\end{proposition}

\begin{proof}
By Proposition~\ref{prop:path-bijection}, the path
$(e_1,e_2,\dots,e_d)$ is a directed path of length $d-1$ from
$x_u$ to $x_v$ in $l(\mathcal{G})$.
Assume for contradiction that there is a shorter path of length
$\ell<d-1$ between $x_u$ and $x_v$ in $\mcal{E}_{\ell}$.
The surjectivity of $\Phi$ then yields a directed path of length
$\ell+1<d$ from $u$ to $v$ in $\mathcal{G}$, contradicting the
minimality of $d$.
Hence the distance in $l(\mathcal{G})$ equals $d-1$.
\end{proof}

\noindent
Propositions~\ref{prop:path-bijection} confirms that there is a one-to-one mapping of reasoning paths between $l(\mcal{G})$ and $\mcal{G}$, therefore we can perform path-based learning and inference in the line graph. Propositions~\ref{prop:length-reduction} implies that a $K$-layer GNN model in $l(\mcal{G})$ is equivalent in the receptive field as a $(K+1)$-layer GNN model in graph $G$, therefore the receptive field of a GNN model increases via line graph transformation.

\section{Datasets}
\label{dataset}
WebQSP is a benchmark dataset for KGQA, derived from the original WebQuestions dataset~\cite{berant2013semantic}. It comprises 4,737 natural language questions annotated with full semantic parses in the form of SPARQL queries executable against Freebase. The dataset emphasizes single-hop questions, typically involving a direct relation between the question and answer entities.

CWQ dataset extends the WebQSP dataset to address more challenging multi-hop question answering scenarios. It contains 34,689 complex questions that require reasoning over multiple facts and relations. Each question is paired with a SPARQL query and corresponding answers, facilitating evaluation in both semantic parsing and information retrieval contexts. The datasets statistics can be found in Table~\ref{tab:dataset_stats}. 

\begin{table}[!t]
\centering
\caption{Dataset statistics and distribution of answer set sizes.}
\label{tab:dataset_stats}
\begin{adjustbox}{width=0.9\textwidth}
\begin{tabular}{c|cc|cccc}
\toprule
\multirow{2}{*}{Dataset} & \multicolumn{2}{c|}{Dataset Size} & \multicolumn{4}{c}{Distribution of Answer Set Size} \\
\cmidrule{2-7}
 & \#Train & \#Test & $\#\text{Ans}=1$ & $2 \leq \#\text{Ans} \leq 4$ & $5 \leq \#\text{Ans} \leq 9$ & $\#\text{Ans} \geq 10$ \\
\midrule
WebQSP & 2,826 & 1,628 & 51.2\% & 27.4\% & 8.3\% & 12.1\% \\
CWQ    & 27,639 & 3,531 & 70.6\% & 19.4\% & 6.0\% & 4.0\% \\
\bottomrule
\end{tabular}
\end{adjustbox}
\end{table}


Following previous practice, we adopt the same training and test split, with the same subgraph construction for each question-answer pair to ensure fairness~\cite{jiang2022unikgqa,luo2024rog,li2024simple,mavromatis2024gnn}.

\section{More Details on Experimental Setup and Implementations}
\label{app:exp}
\tb{Setup.} For model training, we employ two 2-layer GCNs to enable bidirectional message passing. Each GCN has a hidden dimension of 512. We use the Adam optimizer~\cite{kingma2014adam} with a learning rate of $1\times10^{-3}$, and a batch size of 10. Batch normalization~\cite{ioffe2015batch} is not used, as we observe gradient instability when it is applied. The graph retriever is trained for 15 epochs on both datasets, and model selection is performed using cross-validation based on the validation loss. A dropout rate of 0.2~\cite{srivastava2014dropout} is applied for regularization. When multiple valid paths are available, we randomly sample one as the ground-truth supervision signal at each training step.

For evaluating KGQA performance on the test sets, we first generate answers using the downstream reasoner, followed by answer verification using GPT-4o-mini. To prevent potential information leakage, we decouple the answer generation and verification processes, avoiding inclusion of ground-truth answers in the generation prompts. For answer verification, we adopt chain-of-thought prompting~\cite{wei2022chain} to ensure accurate estimation of Macro-F1 and Hit metrics.

\tb{Implementation.} We utilize \ti{networkx}\cite{hagberg2008exploring} for performing line graph transformations and explore all paths between question entities (source nodes) and answer entities (target nodes), and GPT-4o-mini is used during preprocessing for relation targeting. Our remaining implementations are based on PyTorch\cite{paszke2019pytorchimperativestylehighperformance} and PyTorch Geometric~\cite{fey2019fast}.

\section{Efficiency Analysis}
We evaluate the efficiency of the proposed method and baselines based on three metrics: average runtime, average $\#$ LLM calls, and average $\#$ retrieved triples. As shown in Table~\ref{tab:efficiency}, agentic RAG methods (e.g., ToG) incur significantly higher latency and computational cost due to repeated LLM invocations. Among KG-based RAG methods, approaches employing LLM-based retrievers generally require more time than graph-based retrievers due to the LLM inference. Although GNN-RAG and SubgraphRAG exhibit comparable runtime and LLM calls to \proj, our method is more effective thanks to its design choices specifically tailored for KGQA. Furthermore, \proj{} retrieves no more than 50 triples, benefiting from the path-based inference approach, which allows \proj{} to learn when to stop reasoning, thereby avoiding the need to recall a fixed number of triples  as in SubgraphRAG. As a result, \proj{} enables more efficient downstream reasoning with reduced computes, balancing effiency and effectiveness. 
\begin{table}[!t]
\centering
\caption{Efficiency analysis of different methods on WebQSP dataset.}
\label{tab:efficiency}
\scalebox{0.9}{
\begin{tabular}{@{}c|c|c|c|c@{}}

\toprule
\textbf{Methods}                 & \textbf{Hit}               & \textbf{Avg. Runtime (s)} & \textbf{Avg. \# LLM Calls} & \textbf{Avg. \# Triples} \\ \midrule
RoG     & 85.6 & 8.65  & 2    & 49  \\
ToG     & 75.1 & 19.03 & 13.2 & 410 \\
GNN-RAG & 85.7 & 1.82  & 1    & 27  \\
\multicolumn{1}{l|}{SubgraphRAG} & \multicolumn{1}{l|}{90.1} & 2.63                      & 1                          & 100                      \\ \midrule
Ours    & 92.0 & 2.16  & 1    & 32  \\ \bottomrule
\end{tabular}}
\end{table}

\section{Motivating Examples on Rational Paths}
\label{app:label-rationale}
 In this section, we provide 6 intuitive examples of the claim for each case: \tb{(i)} Multiple shortest paths may exist for a given question, not all of which are semantically meaningful, and \tb{(ii)} Some causally-grounded paths may not be the shortest. Figure~\ref{fig:type1_no1}-\ref{fig:type1_no6} demonstrates supports for case 1, and Figure~\ref{fig:type2_no1}-\ref{fig:type2_no6} demonstrates supports for case 2. 

\begin{figure}[t]
\begin{tcolorbox}[colback=gray!5!white,
                  colframe=gray!75!black,
                  title=WebQTest-923\_e3a2d3d50bac69d563de83a7f72eafc0]
\scriptsize


\textbf{Question:}\\
Which country with religious organization leadership \emph{Noddfa, Treorchy} borders England?

\rule{\linewidth}{0.4pt}

\textbf{Candidate shortest paths:}\\[2pt]
{\raggedright   % <‑‑ add this opening brace + directive
\texttt{\textbf{England\patharrow location.location.adjoin\_s\patharrow m.04dgsfb\patharrow
location.adjoining\_relationship.adjoins\patharrow Wales}}\hfill
\textcolor{red}{(rational)}\\[6pt]

\texttt{England\patharrow law.court\_jurisdiction\_area.courts\patharrow
National Industrial Relations Court\patharrow law.court.jurisdiction\patharrow Wales}\\[6pt]

\texttt{England\patharrow organization.organization\_scope.organizations\_with\_this\_scope\patharrow
Police Federation of England and Wales\patharrow
organization.organization.geographic\_scope\patharrow Wales}\\[6pt]

\texttt{England\patharrow organization.organization\_scope.organizations\_with\_this\_scope\patharrow
BES Utilities\patharrow organization.organization.geographic\_scope\patharrow Wales}\\[6pt]

\ldots % remaining paths
\par}

\rule{\linewidth}{0.4pt}

\textbf{Explanation:}\\[2pt]
{\raggedright
\textbf{The first path:} It encodes a direct geographical‐adjacency relation (\texttt{location.location.adjoin\_s}
followed by \texttt{location.adjoining\_relationship.adjoins}), so it
correctly captures that Wales \emph{borders} England.\par\vspace{6pt}

\textbf{The second path:} It links England and Wales through a shared court
system, reflecting \emph{legal jurisdiction} rather than physical
contiguity; therefore it is not a rational answer.\par\vspace{6pt}

\textbf{The third path:} It relies on an organisation (Police Federation of
England and Wales) that operates in both regions.  Operational scope
signals administrative overlap, not territorial borders.\par\vspace{6pt}

\textbf{The fourth path:} Like the third, it uses an organisation’s
geographic scope (BES Utilities) to connect the two regions, so it conveys
no information about adjacency and is likewise non‑rational.\par}

\end{tcolorbox}
\caption{Motivating example on not all shortest paths are rational paths.}
\label{fig:type1_no1}
\end{figure}

\begin{figure}[t]
\begin{tcolorbox}[colback=gray!5!white,
                  colframe=gray!75!black,
                  title=WebQTest-415\_b6ad66a3f1f515d0688c346e16d202e6]
\scriptsize

\textbf{Question:}\\
What movie with film character named Mr. Woodson did Tupac star in?

\rule{\linewidth}{0.4pt}

\textbf{Candidate shortest paths:}\\[2pt]
{\raggedright
\texttt{\textbf{Tupac Shakur\patharrow film.actor.film\patharrow m.0jz0c4\patharrow film.performance.film\patharrow Gridlock'd}}\hfill
\textcolor{red}{(rational)}\\[6pt]

\texttt{Tupac Shakur\patharrow music.featured\_artist.recordings\patharrow Out The Moon\patharrow music.recording.releases\patharrow Gridlock'd}\\[6pt]

\texttt{Tupac Shakur\patharrow music.featured\_artist.recordings\patharrow Wanted Dead or Alive \\ \patharrow music.recording.releases\patharrow Gridlock'd}\\[6pt]

\texttt{Tupac Shakur\patharrow music.artist.track\_contributions\patharrow m.0nj8wrw\patharrow music.track\_contribution.track \\ \patharrow Out The Moon\patharrow music.recording.releases\patharrow Gridlock'd}\\[6pt]

\texttt{Tupac Shakur\patharrow film.music\_contributor.film\patharrow Def Jam's How to Be a Player\patharrow film.film.produced\_by \\ \patharrow Russell Simmons\patharrow film.producer.films\_executive\_produced\patharrow Gridlock'd}\\[6pt]

\par}

\rule{\linewidth}{0.4pt}

\textbf{Explanation:}\\[2pt]
{\raggedright
\textbf{The first path:} This path follows the relation \texttt{film.actor.film} from Tupac Shakur to a role entity, and then \texttt{film.performance.film} to the film Gridlock'd. It accurately models the actor–character–film linkage, so it is a rational answer.\par\vspace{6pt}

\textbf{The second path:} It connects Tupac Shakur to the film Gridlock'd via a featured music recording. This reflects musical involvement, not acting or character presence, so it does not address the question.\par\vspace{6pt}

\textbf{The third path:} Similar to the second, it identifies Tupac as a featured artist on a song associated with the film. However, musical contributions do not imply he played a film character.\par\vspace{6pt}

\textbf{The fourth path:} This path involves nested musical metadata, eventually reaching Gridlock'd via track and recording releases. It does not establish that Tupac portrayed a character in the movie, hence it is not rational.\par\vspace{6pt}

\textbf{The fifth path:} It connects Tupac to Gridlock'd via his contribution to another film (Def Jam's How to Be a Player), which shares a producer with Gridlock'd. This is an indirect production-based relation, not evidence of his acting in Gridlock'd.\par}
\end{tcolorbox}
\caption{Motivating example on not all shortest paths are rational paths.}
\label{fig:type1_no2}
\end{figure}

\begin{figure}[t]
\begin{tcolorbox}[colback=gray!5!white,
                  colframe=gray!75!black,
                  title=WebQTrn-3696\_b874dcb19fa3a6c4e6037dc13f1f3bc4]
\scriptsize

\textbf{Question:}\\
Which state senator from Georgia took this position at the earliest date?

\rule{\linewidth}{0.4pt}

\textbf{Candidate shortest paths:}\\[2pt]
{\raggedright
\texttt{\textbf{Georgia\patharrow government.political\_district.representatives\patharrow m.030qq3n \\ \patharrow government.government\_position\_held.office\_holder\patharrow Saxby Chambliss}}\hfill
\textcolor{red}{(rational)}\\[6pt]

\texttt{Georgia\patharrow government.political\_district.elections\patharrow United States Senate election in Georgia, 2008\patharrow common.topic.image\patharrow Saxby Chambliss}\\[6pt]

\par}

\rule{\linewidth}{0.4pt}

\textbf{Explanation:}\\[2pt]
{\raggedright
\textbf{The first path:} This path directly follows \texttt{government.political\_district.representatives} from Georgia to a position held by an entity, and then uses \texttt{government.government\_position\_held.office\_holder} to identify the person (Saxby Chambliss) who held the position. This correctly models a government role held by someone representing the district, making it a rational and temporally grounded path for determining who assumed the position earliest.\par\vspace{6pt}

\textbf{The second path:} This path connects Georgia to a 2008 Senate election, and then to Saxby Chambliss via an image relation. Although the entity appears in the context of the election, the path does not encode any formal position-holding information or temporal precedence. It is therefore unrelated to identifying the office-holder or the date of assuming the position, and is non-rational.\par}

\end{tcolorbox}
\caption{Motivating example on not all shortest paths are rational paths.}
\label{fig:type1_no3}
\end{figure}


\begin{figure}[t]
\begin{tcolorbox}[colback=gray!5!white,
                  colframe=gray!75!black,
                  title=WebQTrn-3763\_c707414f103503f2530fc654a85645fe]
\scriptsize

\textbf{Question:}\\
What country close to Russia has a religious organization named Ukrainian Greek Catholic Church?

\rule{\linewidth}{0.4pt}

\textbf{Candidate shortest paths:}\\[2pt]
{\raggedright
\texttt{\textbf{Ukrainian Greek Catholic Church\patharrow religion.religious\_organization.leaders\patharrow m.05tnwqd \\ \patharrow religion.religious\_organization\_leadership.jurisdiction\patharrow Ukraine}}\hfill
\textcolor{red}{(rational)}\\[6pt]

\texttt{Russia\patharrow location.location.partially\_contains\patharrow Seym River\patharrow geography.river.basin\_countries\patharrow Ukraine}\\[6pt]

\texttt{Russia\patharrow olympics.olympic\_participating\_country.olympics\_participated\_in\patharrow 2010 Winter Olympics \\ \patharrow olympics.olympic\_games.participating\_countries\patharrow Ukraine}\\[6pt]

\texttt{Russia\patharrow organization.organization\_founder.organizations\_founded\patharrow Commonwealth of Independent States \\ \patharrow organization.organization.founders\patharrow Ukraine}\\[6pt]

\texttt{Russia\patharrow location.location.adjoin\_s\patharrow m.02wj9d3\patharrow location.adjoining\_relationship.adjoins\patharrow Ukraine}\\[6pt]

\par}

\rule{\linewidth}{0.4pt}

\textbf{Explanation:}\\[2pt]
{\raggedright
\textbf{The first path:} This path connects the Ukrainian Greek Catholic Church via \texttt{religion.religious\_organization.leaders} to a leadership entity, then via \texttt{religion.religious\_organization\_leadership.jurisdiction} to Ukraine. It correctly encodes the organizational jurisdiction of the church and identifies the relevant country, making it a rational answer.\par\vspace{6pt}

\textbf{The second path:} This connects Russia to the Seym River and then to Ukraine via river basin membership. It reflects geographic proximity but does not capture any religious organizational structure. Hence, it is non-rational.\par\vspace{6pt}

\textbf{The third path:} This path shows that both Russia and Ukraine participated in the same Olympic games. While it may imply contemporaneity or international context, it says nothing about religious institutions or jurisdictions.\par\vspace{6pt}

\textbf{The fourth path:} This connects Russia and Ukraine through the shared founding of the Commonwealth of Independent States. Although it reflects political cooperation, it does not reveal any information about religious affiliation or structure.\par\vspace{6pt}

\textbf{The fifth path:} This path shows Russia shares a border with Ukraine. It satisfies the "close to Russia" part of the question, but lacks any information about the Ukrainian Greek Catholic Church, making it non-rational.\par}

\end{tcolorbox}
\caption{Motivating example on not all shortest paths are rational paths.}
\label{fig:type1_no4}
\end{figure}


\begin{figure}[t]
\begin{tcolorbox}[colback=gray!5!white,
                  colframe=gray!75!black,
                  title=WebQTrn-3548\_c352f5de0efe2369ee74ef2a99973561]
\scriptsize

\textbf{Question:}\\
What city was the birthplace of Charlton Heston and a famous pro athlete who started their career in 2007?

\rule{\linewidth}{0.4pt}

\textbf{Candidate shortest paths:}\\[2pt]
{\raggedright
\texttt{\textbf{Charlton Heston\patharrow people.person.places\_lived\patharrow m.0h28vy2\patharrow people.place\_lived.location\patharrow Los Angeles}}\hfill
\textcolor{red}{(rational)}\\[6pt]

\texttt{Charlton Heston\patharrow people.deceased\_person.place\_of\_death\patharrow Beverly Hills \\ \patharrow base.biblioness.bibs\_location.city\patharrow Los Angeles}\\[6pt]

\texttt{Charlton Heston\patharrow people.person.children\patharrow Fraser Clarke Heston\patharrow people.person.place\_of\_birth \\ \patharrow Los Angeles}\\[6pt]

\ldots % duplicate or extended variants omitted
\par}

\rule{\linewidth}{0.4pt}

\textbf{Explanation:}\\[2pt]
{\raggedright
\textbf{The first path:} This path uses \texttt{people.person.places\_lived} to access a lived-location node and then follows \texttt{people.place\_lived.location} to reach Los Angeles. Since place-of-birth often overlaps with early-life residence, this path is rational in the absence of explicit birth data.\par\vspace{6pt}

\textbf{The second path:} This connects Charlton Heston to Los Angeles via his place of death (Beverly Hills), which is a sub-location of Los Angeles. However, death location is unrelated to birthplace, making this path non-rational.\par\vspace{6pt}

\textbf{The third path:} This path identifies Charlton Heston’s child and retrieves that child’s birthplace (Los Angeles). While it shares a location with the question's answer, it provides no evidence of Charlton Heston’s own birthplace and is therefore non-rational.\par}

\end{tcolorbox}
\caption{Motivating example on not all shortest paths are rational paths.}
\label{fig:type1_no5}
\end{figure}

\begin{figure}[t]
\begin{tcolorbox}[colback=gray!5!white,
                  colframe=gray!75!black,
                  title=WebQTrn-1399\_64fc62dc06d16e612aafb00889d4ada1]
\scriptsize

\textbf{Question:}\\
What is the country close to Russia where Mikheil Saakashvili holds a government position?

\rule{\linewidth}{0.4pt}

\textbf{Candidate shortest paths:}\\[2pt]
{\raggedright
\texttt{\textbf{Mikheil Saakashvili\patharrow government.politician.government\_positions\_held\patharrow m.0j6t55g \\ \patharrow government.government\_position\_held.jurisdiction\_of\_office\patharrow Georgia}}\hfill
\textcolor{red}{(rational)}\\[6pt]

\texttt{Mikheil Saakashvili\patharrow organization.organization\_founder.organizations\_founded \\ \patharrow United National Movement\patharrow organization.organization.geographic\_scope\patharrow Georgia}\\[6pt]

\texttt{Russia\patharrow location.location.partially\_contains\patharrow Diklosmta \\ \patharrow location.location.partially\_containedby\patharrow Georgia}\\[6pt]

\texttt{Russia\patharrow location.country.languages\_spoken\patharrow Osetin Language \\ \patharrow language.human\_language.main\_country\patharrow Georgia}\\[6pt]

\par}

\rule{\linewidth}{0.4pt}

\textbf{Explanation:}\\[2pt]
{\raggedright
\textbf{The first path:} This path follows \texttt{government.politician.government\_positions\_held} to retrieve a position node, and then uses \texttt{government.government\_position\_held.jurisdiction\_of\_office} to reach Georgia. This explicitly identifies the country where Mikheil Saakashvili held a government role, satisfying both the political and geographical parts of the question. It is thus rational.\par\vspace{6pt}

\textbf{The second path:} This path captures that Saakashvili founded an organization with activity in Georgia. While it indicates political involvement, it does not assert that he held a formal government position in the country. Hence, non-rational.\par\vspace{6pt}

\textbf{The third path:} This connects Russia to Georgia via a geographic relation involving Diklosmta. It reflects proximity, but does not involve Saakashvili or government roles—so it is irrelevant to the question.\par\vspace{6pt}

\textbf{The fourth path:} This path links Russia and Georgia through a shared spoken language (Osetin). While this indicates cultural or linguistic ties, it provides no information about Saakashvili’s political role. It is non-rational.\par\vspace{6pt}
\par}

\end{tcolorbox}
\caption{Motivating example on not all shortest paths are rational paths.}
\label{fig:type1_no6}
\end{figure}


\begin{figure}[t]
\begin{tcolorbox}[colback=gray!5!white,
                  colframe=gray!75!black,
                  title=WebQTrn-2946\_93c6dae3d218dbe112d4120f45c93298]
\scriptsize

\textbf{Question:}\\
What team with mascot named Champ did Tyson Chandler play for?

\rule{\linewidth}{0.4pt}

\textbf{Candidate shortest paths:}\\[2pt]
{\raggedright
\texttt{\textbf{Tyson Chandler\patharrow sports.pro\_athlete.teams\patharrow m.0j2jj7v\patharrow sports.sports\_team\_roster.team \\ \patharrow Dallas Mavericks}}\hfill
\textcolor{red}{(rational)}\\[6pt]

\texttt{Tyson Chandler\patharrow sports.pro\_athlete.teams\patharrow m.0110h779\patharrow sports.sports\_team\_roster.team \\ \patharrow Dallas Mavericks}\\[6pt]

\texttt{Champ\patharrow sports.mascot.team\patharrow Dallas Mavericks}\hfill
\textcolor{blue}{(shortest path)}\\[6pt]
\par}

\rule{\linewidth}{0.4pt}

\textbf{Explanation:}\\[2pt]
{\raggedright
\textbf{The first path:} This path uses \texttt{sports.pro\_athlete.teams} to retrieve a team membership record, and \texttt{sports.sports\_team\_roster.team} to reach the team (Dallas Mavericks). It directly connects Tyson Chandler to the team he played for, making it a rational path.\par\vspace{6pt}

\textbf{The second path:} This is structurally identical to the first but uses a different team-roster node. It also reaches Dallas Mavericks through the correct relation pair, so it is equally valid and rational in form—though redundant if the goal is to find \emph{a} team Chandler played for with mascot Champ.\par\vspace{6pt}

\textbf{The third path:} This connects the mascot Champ to the Dallas Mavericks. It identifies the team correctly but does not involve Tyson Chandler. Therefore, on its own, it does not answer the question. It is necessary context but not sufficient, and thus non-rational as a standalone path.\par}

\end{tcolorbox}
\caption{Motivating example on shortest paths may not be rational paths.}
\label{fig:type2_no1}
\end{figure}


\begin{figure}[t]
\begin{tcolorbox}[colback=gray!5!white,
                  colframe=gray!75!black,
                  title=WebQTrn-934\_02aae167a8fa9f7d45daab265ac650cd]
\scriptsize

\textbf{Question:}\\
Who held their governmental position from 1786 and was the British General of the Revolutionary War?

\rule{\linewidth}{0.4pt}

\textbf{Candidate shortest paths:}\\[2pt]
{\raggedright
\texttt{\textbf{Kingdom of Great Britain\patharrow military.military\_combatant.military\_commanders\patharrow m.04fttv1 \\ \patharrow military.military\_command.military\_commander\patharrow Charles Cornwallis, 1st Marquess Cornwallis}}\hfill
\textcolor{red}{(rational)}\\[6pt]

\texttt{American Revolutionary War\patharrow base.culturalevent.event.entity\_involved \\ \patharrow Charles Cornwallis, 1st Marquess Cornwallis}\hfill
\textcolor{blue}{(shortest)}
\\[6pt]
\ldots
\par}

\rule{\linewidth}{0.4pt}

\textbf{Explanation:}\\[2pt]
{\raggedright
\textbf{The first path:} This path begins with the Kingdom of Great Britain, follows \texttt{military.military\_combatant.military\_commanders} to a command structure, and then \texttt{military.military\_command.military\_commander} to Charles Cornwallis. It directly encodes his role as a British military commander, making it a rational path aligned with the question.\par\vspace{6pt}

\textbf{The second path:} This connects Charles Cornwallis to the American Revolutionary War via an \texttt{entity\_involved} relation. While it establishes that he was involved in the war, it does not specify a command role nor relate to the governmental position, so it is non-rational.
\par}

\end{tcolorbox}
\caption{Motivating example on shortest paths may not be rational paths.}
\label{fig:type2_no2}
\end{figure}


\begin{figure}[t]
\begin{tcolorbox}[colback=gray!5!white,
                  colframe=gray!75!black,
                  title=WebQTrn-2349\_e831da3802943dad506eb1e3fb611847]
\scriptsize

\textbf{Question:}\\
What are the official bird and flower of the state whose capital is Lansing?

\rule{\linewidth}{0.4pt}

\textbf{Candidate shortest paths:}\\[2pt]
{\raggedright
\texttt{\textbf{Lansing\patharrow base.biblioness.bibs\_location.state\patharrow Michigan \\ \patharrow government.governmental\_jurisdiction.official\_symbols\patharrow m.04st85s \\ \patharrow location.location\_symbol\_relationship.symbol\patharrow American robin}}\hfill
\textcolor{red}{(rational)}\\[6pt]


\texttt{State bird\patharrow location.offical\_symbol\_variety.symbols\_of\_this\_kind\patharrow m.04st83j \\ \patharrow location.location\_symbol\_relationship.symbol\patharrow American robin}\hfill
\textcolor{blue}{(shortest)}\\[6pt]

\texttt{State flower\patharrow location.offical\_symbol\_variety.symbols\_of\_this\_kind\patharrow m.0hz8zmz \\ \patharrow location.location\_symbol\_relationship.symbol\patharrow Apple Blossom}\hfill
\textcolor{blue}{(shortest)}\\[6pt]
\ldots
\par}

\rule{\linewidth}{0.4pt}

\textbf{Explanation:}\\[2pt]
{\raggedright
\textbf{The first path:} Begin at Lansing, proceed through \texttt{base.biblioness.bibs\_location.state} to reach Michigan, and then use \texttt{government.governmental\_jurisdiction.official\_symbols} to retrieve the state’s official bird and flower via \texttt{location.location\_symbol\_relationship.symbol}. It correctly model the semantic intent of the question and are rational.\par\vspace{6pt}

\textbf{The second path:} This begins from the general category “State bird” and navigates to the American robin, but it lacks a connection to Michigan or Lansing, so it cannot determine the relevant state’s identity. Thus, it is non-rational.\par\vspace{6pt}

\textbf{The third path:} Similar to the third, this links the symbolic category “State flower” to Apple Blossom, but it does not connect this symbol to any particular state. It is non-rational on its own.
\par}

\end{tcolorbox}
\caption{Motivating example on shortest paths may not be rational paths.}
\label{fig:type2_no3}
\end{figure}


\begin{figure}[t]
\begin{tcolorbox}[colback=gray!5!white,
                  colframe=gray!75!black,
                  title=WebQTrn-88\_54f1262a1dbcb7b82f5b8ebd614401b9]
\scriptsize

\textbf{Question:}\\
In what city and state is the university that publishes the newspaper titled \emph{Santa Clara}?

\rule{\linewidth}{0.4pt}

\textbf{Candidate shortest paths:}\\[2pt]
{\raggedright
\texttt{\textbf{Santa Clara\patharrow education.school\_newspaper.school\patharrow Santa Clara University\patharrow \\ location.location.containedby\patharrow Santa Clara}}\hfill
\textcolor{red}{(rational)}\\[6pt]

\texttt{\textbf{Santa Clara\patharrow education.school\_newspaper.school\patharrow Santa Clara University\patharrow \\ location.location.containedby\patharrow California}}\hfill
\textcolor{red}{(rational)}\\[6pt]

\texttt{Santa Clara\patharrow book.newspaper.circulation\_areas\patharrow California}\hfill
\textcolor{blue}{(shortest)}\\[6pt]

\texttt{Santa Clara\patharrow book.newspaper.headquarters\patharrow Santa Clara\patharrow \\ location.mailing\_address.state\_province\_region\patharrow California}\\[6pt]
\ldots
\par}

\rule{\linewidth}{0.4pt}

\textbf{Explanation:}\\[2pt]
{\raggedright
\textbf{The first and second paths:} These follow \texttt{education.school\_newspaper.school} to reach Santa Clara University, then use \texttt{location.location.containedby} to identify its city (Santa Clara) and state (California). These directly trace the geographic location of the university that publishes the newspaper, making both paths rational.\par\vspace{6pt}

\textbf{The third path:} This connects the newspaper \emph{Santa Clara} to California via a general \texttt{book.newspaper.circulation\_areas} relation. While it suggests distribution within California, it does not ground the newspaper in a specific institution or location, so it is non-rational.\par\vspace{6pt}

\textbf{The fourth path:} This path uses the newspaper's headquarters to reach a city and then derives the state from a mailing address field. It circumvents the university relation required by the question and therefore does not directly answer it; it is non-rational.\par}
\end{tcolorbox}
\caption{Motivating example on shortest paths may not be rational paths.}
\label{fig:type2_no4}
\end{figure}


\begin{figure}[t]
\begin{tcolorbox}[colback=gray!5!white,
                  colframe=gray!75!black,
                  title=WebQTrn-2591\_9f8fc8341d7c53fe16d94a6a23638ec4]
\scriptsize

\textbf{Question:}\\
What country using the Malagasy Ariary currency is China's trading partner?

\rule{\linewidth}{0.4pt}

\textbf{Candidate shortest paths:}\\[2pt]
{\raggedright
\texttt{\textbf{China\patharrow location.statistical\_region.places\_exported\_to\patharrow m.04bfg2f\patharrow \\  location.imports\_and\_exports.exported\_to\patharrow Madagascar}}\hfill
\textcolor{red}{(rational)}\\[6pt]

\texttt{Malagasy ariary\patharrow finance.currency.countries\_used\patharrow Madagascar}\hfill
\textcolor{blue}{(shortest)}\\[6pt]

\texttt{China\patharrow travel.travel\_destination.tour\_operators\patharrow Bunnik Tours\patharrow \\ travel.tour\_operator.travel\_destinations\patharrow Madagascar}\\[6pt]
\ldots
\par}

\rule{\linewidth}{0.4pt}

\textbf{Explanation:}\\[2pt]
{\raggedright
\textbf{The first path:} This path uses \texttt{location.statistical\_region.places\_exported\_to} followed by \texttt{location.imports\_and\_exports.exported\_to} to connect China to Madagascar through a trade relationship. It correctly captures the fact that Madagascar is one of China's trading partners, making it rational.\par\vspace{6pt}

\textbf{The second path:} This identifies Madagascar as a country that uses the Malagasy Ariary, but it does not involve China or any trade relationship. It is relevant as supporting context for currency, but not sufficient to answer the question on its own.\par\vspace{6pt}

\textbf{The third path:} This connects China to Madagascar via a tourism relationship involving a tour operator (Bunnik Tours). While it shows interaction between the countries, it does not concern trade and is thus non-rational in the context of the question.\par}

\end{tcolorbox}
\caption{Motivating example on shortest paths may not be rational paths.}
\label{fig:type2_no5}
\end{figure}


\begin{figure}[t]
\begin{tcolorbox}[colback=gray!5!white,
                  colframe=gray!75!black,
                  title=WebQTrn-2237\_723ad981dc68e3cfe82e7134c8ca8fdb]
\scriptsize

\textbf{Question:}\\
Where are some places to stay at in the city where Gavin Newsom is a government officer?

\rule{\linewidth}{0.4pt}

\textbf{Candidate shortest paths:}\\[2pt]
{\raggedright
\texttt{\textbf{Gavin Newsom\patharrow government.political\_appointer.appointees\patharrow m.03k0n\_d\patharrow \\ government.government\_position\_held.jurisdiction\_of\_office\patharrow San Francisco\patharrow \\ travel.travel\_destination.accommodation\patharrow W San Francisco}}\hfill
\textcolor{red}{(rational)}\\[6pt]

\texttt{Gavin Newsom\patharrow people.person.place\_of\_birth\patharrow San Francisco\patharrow \\ travel.travel\_destination.accommodation\patharrow W San Francisco}\hfill
\textcolor{blue}{(shortest)}\\[6pt]

\texttt{Gavin Newsom\patharrow people.person.place\_of\_birth\patharrow San Francisco\patharrow \\ travel.travel\_destination.accommodation\patharrow Hostelling International, City Center}\hfill
\textcolor{blue}{(shortest)}\\[6pt]
\ldots
\par}

\rule{\linewidth}{0.4pt}

\textbf{Explanation:}\\[2pt]
{\raggedright
\textbf{The first path:} This path links Gavin Newsom to San Francisco through a government appointee role and then uses \texttt{government.government\_position\_held.jurisdiction\_of\_office} to specify the city in which he held office. From there, it identifies accommodations such as W San Francisco via \texttt{travel.travel\_destination.accommodation}. It directly answers the question and is rational.\par\vspace{6pt}

\textbf{The second and third paths:} These use \texttt{people.person.place\_of\_birth} to link Gavin Newsom to San Francisco and then list accommodations in that city. While San Francisco is both his birthplace and his jurisdiction of office, these paths do not establish that the city is where he was a government officer. Thus, they are non-rational according to the expected relation chain.\par}
\end{tcolorbox}
\caption{Motivating example on shortest paths may not be rational paths.}
\label{fig:type2_no6}
\end{figure}


\section{Demonstrations on Retrieved Reasoning Paths from \proj}
In this section, we provide 10 examples to demonstrate the qualitative results of \proj, as shown in Figure~\ref{fig:example1}-\ref{fig:example10}.

\begin{figure}[t]
\begin{tcolorbox}[colback=gray!5!white,
                  colframe=gray!75!black,
                  title=WEBQSP-WebQTest-7]
\scriptsize

\textbf{Question:}\\
Where was George Washington Carver from?

\rule{\linewidth}{0.4pt}

\textbf{Retrieved Paths:}\\[2pt]
{\raggedright
\texttt{George Washington Carver\patharrow people.person.nationality\patharrow United States of America}\\[6pt]

\texttt{George Washington Carver\patharrow people.deceased\_person.place\_of\_death\patharrow Tuskegee}\\[6pt]

\texttt{George Washington Carver\patharrow people.person.places\_lived\patharrow m.03prs0h}\\[6pt]

\texttt{\textbf{George Washington Carver\patharrow people.person.place\_of\_birth\patharrow Diamond}}\\[6pt]

\texttt{George Washington Carver\patharrow people.person.places\_lived\patharrow m.03ppx0s}\\[6pt]

\texttt{George Washington Carver\patharrow people.person.education\patharrow m.04hdfv4}\\[6pt]

\texttt{George Washington Carver\patharrow people.person.education\patharrow m.04hdfv4\patharrow \\ education.education.institution\patharrow Iowa State University}\\[6pt]

\texttt{George Washington Carver\patharrow people.person.places\_lived\patharrow m.03ppx0s\patharrow \\ people.place\_lived.location\patharrow Tuskegee}\\[6pt]

\texttt{George Washington Carver\patharrow people.person.places\_lived\patharrow m.03prs0h\patharrow \\ people.place\_lived.location\patharrow Joplin}\\[6pt]

\texttt{George Washington Carver\patharrow people.person.place\_of\_birth\patharrow Diamond\patharrow \\ location.statistical\_region.population\patharrow m.0hlfnly}\\[6pt]

\par}

\rule{\linewidth}{0.4pt}

\textbf{Ground-truth:}\\
Diamond

\end{tcolorbox}
\caption{Example on the retrieved reasoning paths by \proj}
\label{fig:example1}
\end{figure}


\begin{figure}[t]
\begin{tcolorbox}[colback=gray!5!white,
                  colframe=gray!75!black,
                  title=WEBQSP-WebQTest-928]
\scriptsize

\textbf{Question:}\\
What colleges did Harper Lee attend?

\rule{\linewidth}{0.4pt}

\textbf{Retrieved Paths:}\\[2pt]
{\raggedright
\texttt{Harper Lee\patharrow people.person.education\patharrow m.0lwxmyl\patharrow \\  education.education.institution\patharrow Monroe County High School}\\[6pt]

\texttt{Harper Lee\patharrow people.person.education\patharrow m.0lwxmy1\patharrow \\  education.education.institution\patharrow Huntingdon College}\\[6pt]

\texttt{Harper Lee\patharrow people.person.education\patharrow m.0lwxmy9\patharrow  \\ education.education.institution\patharrow University of Oxford}\\[6pt]

\texttt{Harper Lee\patharrow people.person.education\patharrow m.0n1l46h\patharrow \\  education.education.institution\patharrow University of Alabama School of Law}\\[6pt]

\texttt{Harper Lee\patharrow people.person.education\patharrow m.04hx138\patharrow \\  education.education.institution\patharrow University of Alabama}\\[6pt]

\par}

\rule{\linewidth}{0.4pt}

\textbf{Ground-truth:}\\
University of Alabama, Huntingdon College, University of Oxford, University of Alabama School of Law

\end{tcolorbox}
\caption{Example on the retrieved reasoning paths by \proj}
\label{fig:example2}
\end{figure}


\begin{figure}[t]
\begin{tcolorbox}[colback=gray!5!white,
                  colframe=gray!75!black,
                  title=WEBQSP-WebQTest-1205]
\scriptsize

\textbf{Question:}\\
Who plays Harley Quinn?

\rule{\linewidth}{0.4pt}

\textbf{Retrieved Paths:}\\[2pt]
{\raggedright
\texttt{Harley Quinn\patharrow tv.tv\_character.appeared\_in\_tv\_program\patharrow m.02wm17r\patharrow \\ tv.regular\_tv\_appearance.actor\patharrow Hynden Walch}\\[6pt]

\texttt{Harley Quinn\patharrow cvg.game\_character.games\patharrow m.09dycc\_\patharrow \\ cvg.game\_performance.voice\_actor\patharrow Arleen Sorkin}\\[6pt]

\texttt{Harley Quinn\patharrow film.film\_character.portrayed\_in\_films\patharrow m.0j6pcwz\patharrow \\ film.performance.actor\patharrow Chrissy Kiehl}\\[6pt]

\texttt{Harley Quinn\patharrow tv.tv\_character.appeared\_in\_tv\_program\patharrow m.02wm18b\patharrow \\ tv.regular\_tv\_appearance.actor\patharrow Mia Sara}\\[6pt]

\texttt{Harley Quinn\patharrow tv.tv\_character.appeared\_in\_tv\_program\patharrow m.0wz39vs\patharrow \\ tv.regular\_tv\_appearance.actor\patharrow Arleen Sorkin}\\[6pt]

\par}

\rule{\linewidth}{0.4pt}

\textbf{Ground-truth:}\\
Mia Sara, Hynden Walch, Arleen Sorkin

\end{tcolorbox}
\caption{Example on the retrieved reasoning paths by \proj}
\label{fig:example3}
\end{figure}


\begin{figure}[t]
\begin{tcolorbox}[colback=gray!5!white,
                  colframe=gray!75!black,
                  title=WEBQSP-WebQTest-1599]
\scriptsize

\textbf{Question:}\\
Where did Kim Jong-il die?

\rule{\linewidth}{0.4pt}

\textbf{Retrieved Paths:}\\[2pt]
{\raggedright
\texttt{Kim Jong-il\patharrow people.deceased\_person.place\_of\_burial\patharrow Kumsusan Palace of the Sun}\\[6pt]

\texttt{\textbf{Kim Jong-il\patharrow people.deceased\_person.place\_of\_death\patharrow Pyongyang}}\\[6pt]

\texttt{Kim Jong-il\patharrow people.deceased\_person.place\_of\_death\patharrow Pyongyang\patharrow \\ periodicals.newspaper\_circulation\_area.newspapers\patharrow Rodong Sinmun}\\[6pt]

\texttt{Kim Jong-il\patharrow people.deceased\_person.place\_of\_death\patharrow Pyongyang\patharrow \\ location.location.contains\patharrow Munsu Water Park}\\[6pt]

\texttt{Kim Jong-il\patharrow people.deceased\_person.place\_of\_death\patharrow Pyongyang\patharrow \\ location.location.contains\patharrow Pyongyang University of Science and Technology}\\[6pt]

\par}

\rule{\linewidth}{0.4pt}

\textbf{Ground-truth:}\\
Pyongyang

\end{tcolorbox}
\caption{Example on the retrieved reasoning paths by \proj}
\label{fig:example4}
\end{figure}


\begin{figure}[t]
\begin{tcolorbox}[colback=gray!5!white,
                  colframe=gray!75!black,
                  title=WEBQSP-WebQTest-1993]
\scriptsize

\textbf{Question:}\\
What language do they speak in Argentina?

\rule{\linewidth}{0.4pt}

\textbf{Retrieved Paths:}\\[2pt]
{\raggedright
\texttt{Argentina\patharrow location.country.languages\_spoken\patharrow Spanish Language}\\[6pt]

\texttt{Argentina\patharrow location.country.languages\_spoken\patharrow Yiddish Language}\\[6pt]

\texttt{Argentina\patharrow location.country.languages\_spoken\patharrow Guaraní language}\\[6pt]

\texttt{Argentina\patharrow location.country.languages\_spoken\patharrow Quechuan languages}\\[6pt]

\texttt{Argentina\patharrow location.country.languages\_spoken\patharrow Italian Language}\\[6pt]

\par}

\rule{\linewidth}{0.4pt}

\textbf{Ground-truth:}\\
Yiddish Language, Spanish Language, Quechuan languages, Italian Language, Guaraní language

\end{tcolorbox}
\caption{Example on the retrieved reasoning paths by \proj}
\label{fig:example5}
\end{figure}


\begin{figure}[t]
\begin{tcolorbox}[colback=gray!5!white,
                  colframe=gray!75!black,
                  title=CWQ-WebQTrn-962\_f0c57985929ee8b823983f6e5f104971]
\scriptsize

\textbf{Question:}\\
What actor played a kid in the film with a character named Veteran at War Rally?

\rule{\linewidth}{0.4pt}

\textbf{Retrieved Paths:}\\[2pt]
{\raggedright
\texttt{Forrest Gump\patharrow film.film\_character.portrayed\_in\_films\patharrow m.0jycvw\patharrow \\ film.performance.actor\patharrow Tom Hanks}\\[6pt]

\texttt{\textbf{Forrest Gump\patharrow film.film\_character.portrayed\_in\_films\patharrow m.02xgww5\patharrow  \\ film.performance.actor\patharrow Michael Connor Humphreys}}\\[6pt]

\texttt{Forrest Gump\patharrow common.topic.notable\_for\patharrow g.1258qx91g}\\[6pt]

\texttt{Veteran at War Rally\patharrow common.topic.notable\_for\patharrow g.12z7tmqks\patharrow \\ film.performance.character\patharrow Veteran at War Rally}\\[6pt]

\texttt{Veteran at War Rally\patharrow film.film\_character.portrayed\_in\_films\patharrow m.0y55311\patharrow \\ film.performance.actor\patharrow Jay Ross}\\[6pt]

\texttt{Veteran at War Rally\patharrow film.film\_character.portrayed\_in\_films\patharrow m.0y55311\patharrow \\ film.performance.character\patharrow Veteran at War Rally}\\[6pt]

\texttt{Forrest Gump\patharrow common.topic.notable\_for\patharrow g.1258qx91g\patharrow \\ book.book\_character.appears\_in\_book\patharrow Forrest Gump}\\[6pt]

\par}

\rule{\linewidth}{0.4pt}

\textbf{Ground-truth:}\\
Michael Connor Humphreys

\end{tcolorbox}
\caption{Example on the retrieved reasoning paths by \proj}
\label{fig:example6}
\end{figure}


\begin{figure}[t]
\begin{tcolorbox}[colback=gray!5!white,
                  colframe=gray!75!black,
                  title=CWQ-WebQTrn-2069\_9a4491f5f6a880a03bd96b8180bace4c]
\scriptsize

\textbf{Question:}\\
If I were to visit the governmental jurisdiction where Ricardo Lagos holds an office, what languages do I need to learn to speak?

\rule{\linewidth}{0.4pt}

\textbf{Retrieved Paths:}\\[2pt]
{\raggedright
\texttt{\textbf{Ricardo Lagos\patharrow government.politician.government\_positions\_held\patharrow m.0nbbvk0\patharrow \\ government.government\_position\_held.office\_position\_or\_title\patharrow President of Chile}}\\[6pt]

\texttt{\textbf{Ricardo Lagos\patharrow people.person.nationality\patharrow Chile\patharrow \\ location.country.official\_language\patharrow Spanish Language}}\\[6pt]

\texttt{Ricardo Lagos\patharrow people.person.nationality\patharrow Chile\patharrow \\ base.mystery.cryptid\_area\_of\_occurrence.cryptid\_s\_found\_here\patharrow Giglioli's Whale}\\[6pt]

\texttt{Ricardo Lagos\patharrow people.person.place\_of\_birth\patharrow Santiago\patharrow \\ location.location.contains\patharrow Torre Santa Maria}\\[6pt]

\texttt{Ricardo Lagos\patharrow people.person.nationality\patharrow Chile\patharrow \\  location.country.second\_level\_divisions\patharrow Choapa Province}\\[6pt]

\texttt{Ricardo Lagos\patharrow people.person.place\_of\_birth\patharrow Santiago\patharrow \\ location.location.time\_zones\patharrow Chile Time Zone}\\[6pt]

\texttt{Ricardo Lagos\patharrow people.person.place\_of\_birth\patharrow Santiago\patharrow \\ location.location.contains\patharrow Santiago Metropolitan Park}\\[6pt]

\par}

\rule{\linewidth}{0.4pt}

\textbf{Ground-truth:}\\
Spanish Language, Mapudungun Language, Aymara language, Rapa Nui Language, Puquina Language

\end{tcolorbox}
\caption{Example on the retrieved reasoning paths by \proj}
\label{fig:example7}
\end{figure}


\begin{figure}[t]
\begin{tcolorbox}[colback=gray!5!white,
                  colframe=gray!75!black,
                  title=CWQ-WebQTrn-105\_64489ea2de4b116070d33a0ebcfd4866]
\scriptsize

\textbf{Question:}\\
What currency is used in the country in which Atef Sedki held office in 2013?

\rule{\linewidth}{0.4pt}

\textbf{Retrieved Paths:}\\[2pt]
{\raggedright
\texttt{\textbf{Atef Sedki\patharrow government.politician.government\_positions\_held\patharrow m.0g9442f\patharrow \\ government.government\_position\_held.jurisdiction\_of\_office\patharrow Egypt}}\\[6pt]

\texttt{Atef Sedki\patharrow people.person.place\_of\_birth\patharrow Tanta\patharrow location.location.people\_born\_here\patharrow Atef Ebeid}\\[6pt]

\texttt{Atef Sedki\patharrow people.person.place\_of\_birth\patharrow Tanta\patharrow \\ location.statistical\_region.population\patharrow g.1jmcbhfn7}\\[6pt]

\texttt{Atef Sedki\patharrow people.person.place\_of\_birth\patharrow Tanta\patharrow \\ location.location.people\_born\_here\patharrow Ramadan Abdel Rehim Mansour}\\[6pt]

\texttt{Atef Sedki\patharrow people.person.nationality\patharrow Egypt\patharrow \\ organization.organization\_scope.organizations\_with\_this\_scope\patharrow Reform and Development Misruna Party}\\[6pt]

\texttt{Atef Sedki\patharrow people.person.nationality\patharrow Egypt\patharrow \\ organization.organization\_scope.organizations\_with\_this\_scope\patharrow Islamist Bloc}\\[6pt]

\texttt{Atef Sedki\patharrow people.person.nationality\patharrow Egypt\patharrow \\ location.location.events\patharrow Crusader invasions of Egypt}\\[6pt]

\texttt{Atef Sedki\patharrow government.politician.government\_positions\_held\patharrow m.0g9442f\patharrow \\ government.government\_position\_held.jurisdiction\_of\_office\patharrow Egypt\patharrow  \\ organization.organization\_scope.organizations\_with\_this\_scope\patharrow Free Egyptians Party}\\[6pt]

\par}

\rule{\linewidth}{0.4pt}

\textbf{Ground-truth:}\\
Egyptian pound

\end{tcolorbox}
\caption{Example on the retrieved reasoning paths by \proj}
\label{fig:example8}
\end{figure}

\begin{figure}[t]
\begin{tcolorbox}[colback=gray!5!white,
                  colframe=gray!75!black,
                  title=CWQ-WebQTest-213\_cbbd86314870b15371b43439eb40587a]
\scriptsize

\textbf{Question:}\\
What celebrities did Scarlett Johansson have romantic relationships with that ended before 2006?

\rule{\linewidth}{0.4pt}

\textbf{Retrieved Paths:}\\[2pt]
{\raggedright
\texttt{Scarlett Johansson\patharrow base.popstra.celebrity.dated\patharrow m.065q6ym\patharrow \\ base.popstra.dated.participant\patharrow Ryan Reynolds}\\[6pt]

\texttt{Scarlett Johansson\patharrow base.popstra.celebrity.dated\patharrow m.065q9sh\patharrow \\ base.popstra.dated.participant\patharrow Josh Hartnett}\\[6pt]

\texttt{Scarlett Johansson\patharrow base.popstra.celebrity.dated\patharrow m.064jrnt\patharrow \\ base.popstra.dated.participant\patharrow Josh Hartnett}\\[6pt]

\texttt{\textbf{Scarlett Johansson\patharrow base.popstra.celebrity.dated\patharrow m.064tt90\patharrow \\ base.popstra.dated.participant\patharrow Justin Timberlake}}\\[6pt]

\texttt{Scarlett Johansson\patharrow base.popstra.celebrity.breakup\patharrow m.064ttdz\patharrow \\ base.popstra.breakup.participant\patharrow Justin Timberlake}\\[6pt]

\texttt{\textbf{Scarlett Johansson\patharrow base.popstra.celebrity.dated\patharrow m.065q1sp\patharrow \\ base.popstra.dated.participant\patharrow Jared Leto}}\\[6pt]

\texttt{Scarlett Johansson\patharrow base.popstra.celebrity.dated\patharrow m.065ppwb\patharrow \\ base.popstra.dated.participant\patharrow Patrick Wilson}\\[6pt]

\texttt{Scarlett Johansson\patharrow base.popstra.celebrity.dated\patharrow m.065pwcr\patharrow \\ base.popstra.dated.participant\patharrow Topher Grace}\\[6pt]

\texttt{Scarlett Johansson\patharrow base.popstra.celebrity.breakup\patharrow m.064fpc6\patharrow \\ base.popstra.breakup.participant\patharrow nm1157013}\\[6pt]

\texttt{Scarlett Johansson\patharrow base.popstra.celebrity.dated\patharrow m.064fp5j\patharrow \\ base.popstra.dated.participant\patharrow nm1157013}\\[6pt]

\texttt{Scarlett Johansson\patharrow people.person.spouse\_s\patharrow m.0ygrd3d\patharrow \\ people.marriage.spouse\patharrow Ryan Reynolds}\\[6pt]

\par}

\rule{\linewidth}{0.4pt}

\textbf{Ground-truth:}\\
Justin Timberlake, Jared Leto

\end{tcolorbox}
\caption{Example on the retrieved reasoning paths by \proj}
\label{fig:example9}
\end{figure}


\begin{figure}[t]
\begin{tcolorbox}[colback=gray!5!white,
                  colframe=gray!75!black,
                  title=CWQ-WebQTrn-2569\_712922724a260d96fea082856cd21d6b]
\scriptsize

\textbf{Question:}\\
What sports facility is home to both the Houston Astros and Houston Hotshots?

\rule{\linewidth}{0.4pt}

\textbf{Retrieved Paths:}\\[2pt]
{\raggedright
\texttt{Houston Rockets\patharrow sports.sports\_team.arena\_stadium\patharrow Toyota Center}\\[6pt]

\texttt{Houston Hotshots\patharrow sports.sports\_team.arena\_stadium\patharrow NRG Arena}\\[6pt]

\texttt{Houston Rockets\patharrow sports.sports\_team.venue\patharrow m.0wz1znd\patharrow \\ sports.team\_venue\_relationship.venue\patharrow Toyota Center}\\[6pt]

\texttt{Houston Hotshots\patharrow sports.sports\_team.venue\patharrow m.0x2dzn8\patharrow \\ sports.team\_venue\_relationship.venue\patharrow Lakewood Church Central Campus}\\[6pt]

\texttt{\textbf{Houston Rockets\patharrow sports.sports\_team.venue\patharrow m.0wz8qf2\patharrow \\ sports.team\_venue\_relationship.venue\patharrow Lakewood Church Central Campus}}\\[6pt]

\texttt{\textbf{Houston Rockets\patharrow sports.sports\_team.arena\_stadium\patharrow Lakewood Church Central Campus}}\\[6pt]

\texttt{Houston Rockets\patharrow sports.sports\_team.location\patharrow Houston\patharrow \\ travel.travel\_destination.tourist\_attractions\patharrow Toyota Center}\\[6pt]

\par}

\rule{\linewidth}{0.4pt}

\textbf{Ground-truth:}\\
Lakewood Church Central Campus

\end{tcolorbox}
\caption{Example on the retrieved reasoning paths by \proj}
\label{fig:example10}
\end{figure}

\section{Prompt Template}
\label{prompt_template}
We provide the prompt template in this section, for rational paths identification, relation targeting, and hallucination detection, as shown in Figure~\ref{fig:get_rational_paths}-\ref{fig:hallucination_detect}.
\begin{figure}[!t]
\begin{tcolorbox}[colback=gray!5!white,
                  colframe=gray!75!black,
                  title=Prompt template for identifying rational paths]
\scriptsize

\tb{Example} \\ [5pt]
Given a question <\texttt{example question}>, the reasoning paths are: \\[5pt]

<\texttt{reasoning paths}> \\ [5pt]

The rational paths are: \\ [3pt]

<\texttt{Rational Paths}> \\ [5pt]

\tb{Explanation} \\ [3pt]

<\texttt{Explanation}> \\ [5pt]

\tb{Task} \\ [5pt]

Now given question <\texttt{question}>, the reasoning paths are: \\ [5pt]

<\texttt{Candidate Paths}> \\ [5pt]

Identify all the rational paths, and list below, with explanations: \\ [5pt]

<\texttt{Rational Paths}> \\ [3pt]

<\texttt{Explanations}> \\ [3pt]

\rule{\linewidth}{0.4pt}
\end{tcolorbox}
\caption{Prompt template for retrieving rational reasoning paths.}
\label{fig:get_rational_paths}
\end{figure}


\begin{figure}[!t]
\begin{tcolorbox}[colback=gray!5!white,
                  colframe=gray!75!black,
                  title=Prompt template for relation targeting]
\scriptsize

\tb{Example} \\ [5pt]
Given a question <\texttt{example question}>, the question entity is:<\texttt{question entity}>, the candidate relations for <\texttt{question entity}> are: <\texttt{relations}>. The possible relations for this question are: \\[5pt]

<\texttt{relations}>  \\[5pt]

\tb{Task} \\ [5pt]

Now given question <\texttt{question}>, question entity: <\texttt{question entity}>, relations with the entity: <\texttt{relations}>. List the possible relations for this question. \\ [5pt]

<\texttt{relations}> \\ [5pt]
\end{tcolorbox}
\caption{Prompt template for potential relation targeting.}
\label{fig:get_relation_cand}
\end{figure}


\begin{figure}[!t]
\begin{tcolorbox}[colback=gray!5!white,
                  colframe=gray!75!black,
                  title=Prompt template for hallucination detection]
\scriptsize

\tb{Task} \\ [5pt]
Given the question <\texttt{question}>, the retrieved reasoning paths are: \\[5pt]

<\texttt{reasoning paths}>  \\[5pt]

Now answer this question, and indicate whether you used the information provided above to answer it.

<\texttt{answers}> \\ [5pt]

\#\#\# LLM reasoner:\\ [2pt]
I have used the provided knowledge: <\texttt{Yes|No}>. \\ [3pt]

<\texttt{Explanations}>
\end{tcolorbox}
\caption{Prompt template for hallucination detection.}
\label{fig:hallucination_detect}
\end{figure}

\section{Limitations of \proj}
While \proj{} demonstrates strong performance in KGQA tasks and effectively improves generalization via structured graph retrieval, several limitations remain. 
First, our current approach focuses on enhancing retrieval through structural and relational inductive bias, but does not leverage the complementary strengths of LLM-based retrievers. Designing hybrid retrievers that combine the efficiency of graph-based reasoning with the flexibility and expressiveness of LLMs remains an open challenge. Such integration could potentially yield more robust and scalable KGQA systems.
Second, \proj{} assumes the availability of well-formed KGs and does not address errors or missing entities in the graph itself. Our future work may address the issue of the these aspects.

\section{Software and Hardware}
We conduct all experiments using PyTorch~\cite{paszke2019pytorchimperativestylehighperformance} (v2.1.2) and PyTorch Geometric~\cite{fey2019fast} on Linux servers equipped with NVIDIA A100 GPUs (80GB) and CUDA 12.1.
